% part: intuitionistic-logic
% chapter: propositions-as-types
% section: introduction

\documentclass[../../../include/open-logic-section]{subfiles}

\begin{document}

\olfileid{int}{pty}{int}

\olsection{పరిచయం}

లాంబ్డా కలనశాస్త్రం గణనకు ఒక సరళ నమూనా. చరాల నుంచి నిర్మించిన పదాలపై ఇది పనిచేస్తుంది. $N$, $M$ పదాలైతే, $(NM)$ కూడా పదమే ("$N$ను~$M$కు ప్రయోగించడం"). $N$ పదమైతే $\lambd[x][N]$ ఒక పదం. $N$లో చరం~$x$ ఉంటే, $\lambd[x][N]$లోని సంబంధిత $x$ ఘటనలను $\lambd[x]$ ఆపరేటర్ బంధిస్తుంది; కాబట్టి అవి ఇక స్వేచ్ఛా ఘటనలు కావు. $(\lambd[x][N]M)$ రూపంలోని పదం $\Subst{N}{M}{x}$కు $\beta$-సంకోచితం అవుతుందంటాం (అంటే~$N$లోని ప్రతి స్వేచ్ఛా~$x$ ఘటన స్థానంలో $M$ను ఉంచిన ఫలితం). $N$లోని ఒక ఉపపదాన్ని $\beta$-సంకోచనం చేసి $N'$ను పొందితే, పదం~$N$ పదం~$N'$కు $\beta$-పరివర్తితం అవుతుంది; $N \redone N'$ అని రాస్తాం. $\lambd$-కలనశాస్త్రంలో గణన అంటే $N \redone N_1
\redone N_2 \redone \dots$ అనే వరుస. ఏ ఉపపదమూ $\beta$-సంకోచనం చేయలేని $N_k$ వద్ద ఈ వరుస ముగిస్తే, దాన్ని \emph{సాధారణ} రూపం లేదా~$N$ సాధారణ రూపం అంటాం. $\lambd$-కలనశాస్త్రంలోని గణనను ఇటువంటి వరుసగా భావించండి. సాధారణ రూపం వస్తే గణన ఆగుతుంది; ఆ సాధారణ రూపమే గణన ఫలితం. ఈ సరళ గణన నమూనా ట్యూరింగ్-సంపూర్ణమైనది: ప్రతి పాక్షిక పునరావృత్త ప్రమేయాన్ని $\lambd$-కలనశాస్త్రంలోని ఒక పదంతో గణించవచ్చు.

హాస్కెల్ కర్రీ, విలియం హోవర్డ్ స్వతంత్రంగా సహజ నిగమనానికి, $\lambd$-కలనశాస్త్రానికి మధ్య దగ్గరి పోలికను కనుగొన్నారు. అంతర్బోధగా $\lambd[x][N]$ అనే పదం ఒక ప్రమేయం: $x$ దాని ఆర్గ్యుమెంటు;~$x$ ఇచ్చినప్పుడు విలువను ఎలా గణించాలో $N$ చెబుతుంది. $(\lambd[x][N]M)$ అనే పదం $\lambd[x][N]$ ప్రమేయాన్ని ఆర్గ్యుమెంటు~$M$కు ప్రయోగిస్తుంది; $\beta$-సంకోచనం దాని విలువను గణించే విధానాన్ని చెబుతుంది: $N$లో $x$ స్థానంలో~$M$ను ఉంచాలి (ఫలిత పదం విలువను గణించడం కొనసాగించాలి). ఇప్పుడు ఇవి స్థిర నిర్వచన ప్రాంతం, వ్యాప్తి గల ప్రమేయాలని భావించండి. $\lambd[x][N]$ నిర్వచన ప్రాంతం~$A$, వ్యాప్తి~$B$ అయితే, చరం~$x$ ఆర్గ్యుమెంట్లను~$A$ నుంచి మాత్రమే తీసుకుంటుందని, $N$ విలువలను~$B$లో ఉత్పత్తి చేస్తుందని భావించాలి. అందువల్ల $\lambd[x][N]$ అనేది $A \to B$ ప్రమేయమై, $M \in A$ అయితేనే $(\lambd[x][N]M)$ అర్థవంతం. ఈ సమాచారాన్ని $\lambd$-కలనశాస్త్రానికి చేరిస్తే \emph{రకాలతో కూడిన} $\lambd$-కలనశాస్త్రం వస్తుంది. రకాలతో కూడిన $\lambd$-కలనశాస్త్రంలో ప్రతి $(\lambd[x][N]M)$ వ్యక్తీకరణ అనుమతించబడదు; $\lambd[x][N]$, $M$ల "రకాలు" సరిపోవాలి. అంటే $\lambd[x][N]$ రకం $A \lif B$, $M$ రకం~$A$ అయి ఉండాలి. $x$ రకం~$A$, $N$ రకం~$B$ అయితేనే $\lambd[x][N]$ రకం $A \to B$ అవుతుంది. సాధారణంగా ప్రతి $(M_1M_2)$ వ్యక్తీకరణ అనుమతించబడదు. $M_1$ రకం $A \to B$, $M_2$ రకం~$A$ అయినప్పుడే $(M_1M_2)$ అనుమతించబడుతుంది; అప్పుడు $(M_1M_2)$ రకం~$B$.

ఈ రక నిర్ణయ నియమాలు సహజ నిగమనంలోని \tetoken{నిగమనాలకు}{!!{derivation}s} స్పష్టంగా అనురూపం; రకాలను \tetoken{సూత్రాలు}{!!{formula}s} సూచిస్తాయి, \tetoken{నిగమనాలు}{!!{derivation}s} $\lambd$-పదాలకు అనురూపం. ఉదాహరణకు $!A \lif !B$కు గల \tetoken{ఒక నిగమనం}{!!a{derivation}}, ప్రమేయ రకం $!A \lif !B$ కలిగిన $\lambd[\typeof{x}{!A}][N]$ పదానికి అనురూపం. సహజ నిగమన నియమాలు రక లాంబ్డా కలనశాస్త్రంలోని రక నిర్ణయ నియమాలకు అనురూపం; ఉదాహరణకు,
\begin{prooftree}
  \AxiomC{$\Discharge{!A}{x}$}
  \DeduceC{$!B$}
  \DischargeRule{\Intro\lif}{x}
  \UnaryInfC{$!A \lif !B$}
  \DisplayProof\bottomAlignProof
  \quad\text{అనురూపం}\quad
  \Axiom$x: !A \fCenter N: !B$
  \UnaryInf$ \fCenter \lambd[\typeof{x}{!A}][N] : !A \lif !B$
\end{prooftree}
ఇక్కడ కుడి నియమం అర్థం ఇది: $x$ రకం~$!A$, $N$ రకం~$!B$ అయితే, $\lambd[\typeof{x}{!A}][N]$ రకం~$!A \lif
!B$.

$\Elim\lif$ నియమం ప్రయోగ పదాలకు సంబంధించిన రక నిర్ణయ నియమానికి అనురూపం; అంటే,
\[
  \AxiomC{$!A \lif !B$}
  \AxiomC{$!A$}
  \RightLabel{\Elim\lif}
  \BinaryInfC{$!B$}
  \DisplayProof
  \quad\text{అనురూపం}\quad
  \Axiom$\fCenter M_1: !A \lif !B$
  \Axiom$\fCenter M_2: !A$
  \BinaryInf$ \fCenter (M_1M_2) : !B$
  \DisplayProof
\]
\Intro\lif{} వెంటనే \Elim\lif{} వస్తే, \tetoken{నిగమనాన్ని}{!!{derivation}} సరళీకరించవచ్చు:
\begin{gather*}
  \AxiomC{$\Discharge{!A}{x}$}
  \DeduceC{$!B$}
  \DischargeRule{\Intro\lif}{x}
  \UnaryInfC{$!A \lif !B$}
  \AxiomC{}
  \DeduceC{$!A$}
  \RightLabel{\Elim\lif}
  \BinaryInfC{$!B$}
  \DisplayProof
  \quad\redone\quad
  \AxiomC{}
  \DeduceC{$!A$}
  \DeduceC{$!B$}
  \DisplayProof
\end{gather*}
ఇది లాంబ్డా పదాల $\beta$-సంకోచనానికి అనురూపం:
\[
(\lambd[\typeof{x}{!A}][N]M) \quad\redone\quad
\Subst{N}{M}{x}.
\]
\sourcecorrection{OLTEPTPTYINT-001}{మూలం ఇక్కడ composition పదాలు అని పేర్కొంది; చూపిన $(M_1M_2)$ నియమం ప్రమేయాల సమ్మేళనం కోసం కాదు, ప్రయోగం కోసం. పక్కనున్న వివరణకు అనుగుణంగా ప్రయోగ పదాలు అని అనువదించాం.}

$\land$ నియమాలకు "లబ్ధ" రకాలతో, $\lor$ నియమాలకు "యోగ" రకాలతో ఇలాంటి అనురూపతలు ఉన్నాయి.

సరళ రక లాంబ్డా కలనశాస్త్రంలోని పదాలకు, సహజ నిగమన \tetoken{నిగమనాలకు}{!!{derivation}s} మధ్య ఈ అనురూపతను "కర్రీ--హోవర్డ్", "నిరూపణలు ప్రోగ్రామ్‌లుగా", లేదా "ప్రతిపాదనలు రకాలుగా" అనురూపత అంటారు. \tetoken{సూత్రాలు}{!!{formula}s} (ప్రతిపాదనలు) రకాలకు, నిరూపణలు పదాలకు అనురూపం కావడమే కాకుండా, అనురూపతలను ఇలా సంగ్రహించవచ్చు:
\begin{center}
  \begin{tabular}{c c c}
    తర్కం & ప్రోగ్రామ్ \\
    \hline
    ప్రతిపాదన/\tetoken{సూత్రం}{!!{formula}} & రకం \\
    నిరూపణ & పదం \\
    ఊహ & చరం \\
    విడుదలైన ఊహ & బద్ధ చరం \\
    విడుదల కాని ఊహ & స్వేచ్ఛా చరం \\
    $!A \lif !B$ & ప్రమేయ రకం \\
    $!A \land !B$ & లబ్ధ రకం \\
    $!A \lor !B$ & యోగ రకం \\
    $\lfalse$ & ఖాళీ రకం \\
    \hline
  \end{tabular}
\end{center}

కర్రీ--హోవర్డ్ అనురూపత స్వయంచాలిత నిరూపణ సహాయక పద్ధతులకు, ప్రోగ్రామ్‌ల రక తనిఖీ పద్ధతులకు పునాదులలో ఒకటి. ఎందుకంటే ప్రతిపాదనకు సాక్ష్యమైన నిరూపణను తనిఖీ చేయడం (పైన చేసినట్లుగా), ప్రోగ్రామ్ (పదం)కు ప్రకటించిన రకం ఉందో లేదో తనిఖీ చేయడమే.
\end{document}
